% !PRINTED 547
\weberpage{547}
%<<CONT
und keine anderen hat. Darin liegt der Beweis des folgenden Satzes:

3. \gesperrt{Eine Zahl $\omega$ in $\mathfrak{o}$ ist dann und nur dann nach dem Modul $\pi$ mit einer rationalen Zahl congruent, wenn sie der Bedingung}
\[ \omega^p \equiv \omega \pmod{\pi} \]
\gesperrt{genügt.}

Beachtet man noch, dass die Polynomialcoëfficienten in der $p^{\text{ten}}$ Potenz eines Polynoms alle durch $p$ theilbar sind, mit Ausnahme derer, die zu den $p^{\text{ten}}$ Potenzen der einzelnen Glieder des Polynoms gehören, so ergiebt sich noch folgender Satz, der sich auf ganze Functionen in $\Omega$ von beliebigen Veränderlichen $x, y, \ldots$ bezieht:

4. \gesperrt{Ist $\psi(x,y,\ldots)$ eine ganze Function der Variablen $x, y, \ldots$ mit ganzzahligen Coëfficienten aus $\Omega$, so ist das Bestehen der Congruenz}
\[ [\psi(x,y,\ldots)]^p \equiv \psi(x^p, y^p, \ldots) \pmod{\pi} \]
\gesperrt{die nothwendige und hinreichende Bedingung dafür, dass alle Coëfficienten von $\psi$ mit ganzen rationalen Zahlen nach dem Modul $\pi$ congruent sind.}

\begin{center}
\S.\,151.\\
\gesperrt{Die Dedekind'schen Ideale.}
\end{center}

\gesperrt{Dedekind} gründet in den schon oben erwähnten Arbeiten die Theorie der algebraischen Zahlen auf den Begriff des \gesperrt{Ideals}.

Wir wollen jetzt nachweisen, dass die Theorie der Ideale im Wesen übereinstimmt mit der Theorie der Functionale, indem wir zeigen, wie der Uebergang von der einen zur anderen bewirkt werden kann.

Das System aller ganzen Zahlen eines algebraischen Zahlkörpers $\Omega$ soll, wie oben, mit $\mathfrak{o}$ bezeichnet werden\footnote{Vgl. \gesperrt{Dirichlet-Dedekind}, Vorlesungen über Zahlentheorie im \S.\,167 der dritten, \S.\,177 der vierten Auflage.}. Ein in $\mathfrak{o}$ enthaltenes Zahlensystem $\mathfrak{a}$ wird ein \gesperrt{Ideal} genannt, wenn es den beiden Forderungen genügt:
%CONT>>
